\documentclass[../DoAn.tex]{subfiles}
\begin{document}

\begin{center}
    \Large{\textbf{ABSTRACT}}\\
\end{center}
\vspace{1cm}

Floods are among the most frequent natural disasters in Vietnam, particularly in the central region, where they often disrupt communication and hinder rescue operations. Existing disaster assistance applications provide limited support for real-time location sharing, emergency coordination, and transparent management of post-disaster relief activities. These limitations reduce the effectiveness of rescue efforts and make it difficult to coordinate assistance among citizens, rescuers, benefactors, authorities, and administrators.

To address these challenges, this thesis presents FloodHelper, a disaster assistance platform that integrates real-time location sharing, emergency rescue coordination, and charity campaign management within a unified system. The proposed solution consists of a mobile application for end users and a web-based management system for authorities and administrators. The system adopts a hybrid software architecture combining the Client–Server model and the Publish–Subscribe communication pattern. The core technical approach utilizes the MQTT protocol for communication and location sharing, selected specifically for its lightweight footprint and high resilience in weak network environments. To mitigate intermittent connectivity, the mobile application implements a mechanism to save the victim’s last known location, ensuring rescue teams always have a starting reference point.  In addition, the system provides a charity management module that supports campaign management, donation tracking, relief distribution, and financial reporting to improve transparency throughout the relief process.

The developed system successfully implements the major functional requirements identified during the analysis phase, including user management, friend management, real-time location sharing, emergency rescue coordination, charity campaign management, and administrative functions. The implementation demonstrates the feasibility of integrating these services into a single platform that supports multiple stakeholders involved in disaster response and post-disaster recovery.

The main contribution of this thesis is the development of an integrated disaster assistance platform that combines emergency communication with transparent charity management in a unified architecture. By facilitating timely information sharing, improving coordination among different user groups, and enhancing the transparency of relief activities, the proposed system provides a practical foundation for supporting disaster response and community recovery.

\begin{flushright}
    \begin{tabular}{@{}c@{}}
        Student \\
        \textit{Đăng} \\
        \textit{Vũ Hải Đăng}
    \end{tabular}
\end{flushright}
\end{document}